% Options for packages loaded elsewhere
\PassOptionsToPackage{unicode}{hyperref}
\PassOptionsToPackage{hyphens}{url}
\PassOptionsToPackage{dvipsnames,svgnames,x11names}{xcolor}
%
\documentclass[
  11pt,
  a4paper,
]{ltjsarticle}
\usepackage{amsmath,amssymb}
\usepackage{iftex}
\ifPDFTeX
  \usepackage[T1]{fontenc}
  \usepackage[utf8]{inputenc}
  \usepackage{textcomp} % provide euro and other symbols
\else % if luatex or xetex
  \usepackage{unicode-math} % this also loads fontspec
  \defaultfontfeatures{Scale=MatchLowercase}
  \defaultfontfeatures[\rmfamily]{Ligatures=TeX,Scale=1}
\fi
\usepackage{lmodern}
\ifPDFTeX\else
  % xetex/luatex font selection
\fi
% Use upquote if available, for straight quotes in verbatim environments
\IfFileExists{upquote.sty}{\usepackage{upquote}}{}
\IfFileExists{microtype.sty}{% use microtype if available
  \usepackage[]{microtype}
  \UseMicrotypeSet[protrusion]{basicmath} % disable protrusion for tt fonts
}{}
\makeatletter
\@ifundefined{KOMAClassName}{% if non-KOMA class
  \IfFileExists{parskip.sty}{%
    \usepackage{parskip}
  }{% else
    \setlength{\parindent}{0pt}
    \setlength{\parskip}{6pt plus 2pt minus 1pt}}
}{% if KOMA class
  \KOMAoptions{parskip=half}}
\makeatother
\usepackage{xcolor}
\usepackage[margin=24mm]{geometry}
\setlength{\emergencystretch}{3em} % prevent overfull lines
\providecommand{\tightlist}{%
  \setlength{\itemsep}{0pt}\setlength{\parskip}{0pt}}
\setcounter{secnumdepth}{5}
\ifLuaTeX
  \usepackage{selnolig}  % disable illegal ligatures
\fi
\IfFileExists{bookmark.sty}{\usepackage{bookmark}}{\usepackage{hyperref}}
\IfFileExists{xurl.sty}{\usepackage{xurl}}{} % add URL line breaks if available
\urlstyle{same}
\hypersetup{
  colorlinks=true,
  linkcolor={blue},
  filecolor={Maroon},
  citecolor={Blue},
  urlcolor={blue},
  pdfcreator={LaTeX via pandoc}}

\author{}
\date{}

\begin{document}

\hypertarget{ux8a18ux6cd5}{%
\section{記法}\label{ux8a18ux6cd5}}

この教材では、ベクトルを小文字、行列を大文字で表します。

\[
x\in\mathbb{R}^n,\qquad A\in\mathbb{R}^{m\times n}
\]

\texttt{A} の第 \texttt{i} 行第 \texttt{j} 列を \texttt{a\_\{ij\}}
と書きます。

\hypertarget{ux57faux672cux8a18ux53f7}{%
\subsection{基本記号}\label{ux57faux672cux8a18ux53f7}}

\begin{itemize}
\tightlist
\item
  \texttt{x\^{}T}: 転置
\item
  \texttt{A\^{}\{-1\}}: 逆行列
\item
  \texttt{I}: 単位行列
\item
  \texttt{0}: 零ベクトルまたは零行列
\item
  \texttt{span\{v\_1,...,v\_k\}}: ベクトルの一次結合で生成される空間
\item
  \texttt{ker(A)}: \texttt{Ax=0} を満たす入力の集合
\item
  \texttt{Im(A)} / \texttt{Col(A)}: \texttt{A} の像・列空間
\item
  \texttt{rank(A)}: 像の次元
\item
  \texttt{\textbar{}\textbar{}x\textbar{}\textbar{}\_2}: Euclid ノルム
\item
  \texttt{\textless{}x,y\textgreater{}}: 内積
\end{itemize}

\hypertarget{ux5217ux30d9ux30afux30c8ux30ebux3092ux57faux672cux3068ux3059ux308b}{%
\subsection{列ベクトルを基本とする}\label{ux5217ux30d9ux30afux30c8ux30ebux3092ux57faux672cux3068ux3059ux308b}}

特に断らない限り、ベクトルは列ベクトルです。

\[
x=\begin{bmatrix}x_1\\x_2\\\vdots\\x_n\end{bmatrix}
\]

したがって \texttt{A∈R\^{}\{m×n\}} と \texttt{x∈R\^{}n} に対して

\[
Ax\in\mathbb{R}^m
\]

です。

\hypertarget{ux7b49ux3057ux3044ux3068ux540cux578bux3092ux533aux5225ux3059ux308b}{%
\subsection{「等しい」と「同型」を区別する}\label{ux7b49ux3057ux3044ux3068ux540cux578bux3092ux533aux5225ux3059ux308b}}

抽象的なベクトル空間 \texttt{V} と座標空間 \texttt{R\^{}n}
は、基底を選べば対応づけられますが、概念として同一ではありません。行列は「線形写像そのもの」ではなく、「基底を固定したときの線形写像の表現」です。

この区別は、基底変換や固有値分解を理解するときに重要になります。

\end{document}
